\section{度量}
\label{sec:metric-metrics}

平面上两点相距多远，答案取决于怎样走。直线飞行时用欧氏距离；只能沿东西、南北两个方向分别调整、并以较慢的一个方向为准时，用两个坐标差的最大值；只关心两个位置是否相同时，相同记为$0$，不同记为$1$。本节先比较这三种规则，提炼它们共同遵守的条件，得到度量的定义；再用开球把距离变成第~\ref{chap:topological-spaces}章的拓扑，并讨论距离额外提供的有界性与可分性判据。

\subsection{度量的定义}
\label{sec:metric-definition}

取平面上两点$\bm p=(0,0)$与$\bm q=(3,4)$。三种规则给出的距离分别是
\begin{equation*}
d_2(\bm p,\bm q)=\sqrt{3^2+4^2}=5,\qquad
d_\infty(\bm p,\bm q)=\max\{3,4\}=4,\qquad
d_0(\bm p,\bm q)=1,
\end{equation*}
其中$d_2(\bm x,\bm y)=\bigl((x_1-y_1)^2+(x_2-y_2)^2\bigr)^{1/2}$，$d_\infty(\bm x,\bm y)=\max\{\lvert x_1-y_1\rvert,\lvert x_2-y_2\rvert\}$，$d_0(\bm x,\bm y)$在$\bm x=\bm y$时为$0$、否则为$1$。三个数不同，判断“谁离$\bm p$近”的结果也不同：到$\bm p$的距离小于$1$的点，在$d_2$下组成圆盘，在$d_\infty$下组成边长为$2$的正方形，在$d_0$下只有$\bm p$本身。

数值不同，三种规则却有共同之处。每个距离都不小于零，并且只有两点相同时才为零；从$\bm p$到$\bm q$与从$\bm q$到$\bm p$的距离相同；绕经第三点不会比直接比较更近。把这三条作为公理，就得到不依赖坐标的距离概念。

\begin{definition}[度量]
\label{def:metric-metric}
设$X$是非空集合。映射$d:X\times X\to[0,\infty)$若对任意$x,y,z\in X$满足
\begin{enumerate}
\item 零距离判定：$d(x,y)=0$当且仅当$x=y$；
\item 对称性：$d(x,y)=d(y,x)$；
\item 三角不等式：$d(x,z)\leq d(x,y)+d(y,z)$，
\end{enumerate}
则称$d$为$X$上的一个\term{度量}（metric），称$(X,d)$为\term{度量空间}（metric space）。
\end{definition}

非负性已经写在取值范围$[0,\infty)$里。零距离判定保证距离能区分不同的点；对称性说明距离不带方向；三角不等式说明绕道不会缩短距离，它是后面几乎所有估计的来源。度量空间由集合与度量共同确定：同一个$\mathbb R^2$配上$d_2$、$d_\infty$或$d_0$，是三个不同的度量空间，下文提到“空间$X$”时总默认已经选定了度量。

三角不等式还有一个常用的变形。对任意$x,y,z$，由$d(x,z)\leq d(x,y)+d(y,z)$得$d(x,z)-d(y,z)\leq d(x,y)$，交换$x,y$的角色又得$d(y,z)-d(x,z)\leq d(x,y)$，所以
\begin{equation}
\lvert d(x,z)-d(y,z)\rvert\leq d(x,y).
\label{eq:metric-reverse-triangle}
\end{equation}
它说明：$x$移动多少，$x$到固定点$z$的距离至多变化多少。第~\ref{chap:metric-space-maps}章用它证明距离函数连续。

\begin{example}[三种平面距离的核对]
\label{ex:metric-three-metrics}
对$d_\infty$，前两条公理直接成立。三角不等式逐坐标核对：$\lvert x_i-z_i\rvert\leq\lvert x_i-y_i\rvert+\lvert y_i-z_i\rvert\leq d_\infty(\bm x,\bm y)+d_\infty(\bm y,\bm z)$对$i=1,2$都成立，再对$i$取最大值即可。对$d_0$，若$\bm x\neq\bm z$，则$\bm y$至少与其中一点不同，右端至少为$1$；若$\bm x=\bm z$，左端为$0$。欧氏距离$d_2(\bm x,\bm y)=\lVert\bm x-\bm y\rVert_2$的三角不等式就是欧氏范数的Minkowski不等式。更一般地，任意向量空间上的范数$\lVert\cdot\rVert$都给出度量$d(\bm x,\bm y)=\lVert\bm x-\bm y\rVert$；有限维情形的范数及其三角不等式可参考\BookChapterRef{chap:normed-inner-product-spaces}{卷二“赋范空间与内积空间”章}，一般赋范空间在第~\ref{chap:normed-spaces}章讨论。
\end{example}

度量的对象不必是数或向量。设$X=\{0,1\}^n$是长度为$n$的二进制串，$d_H(x,y)$记$x,y$取值不同的位置个数，称为\term{Hamming距离}（Hamming distance）。例如$d_H(0110,0011)=2$。若$x,z$在某个位置不同，则$y$在这个位置至少与$x,z$之一不同，所以$x,z$的每个不同位置都至少计入$d_H(x,y)$与$d_H(y,z)$中的一个，三角不等式成立。

并不是每个看起来合理的“差异量”都是度量。在$\mathbb R$上取$e(x,y)=(x-y)^2$，前两条公理成立，但$e(0,2)=4$，而$e(0,1)+e(1,2)=2$，三角不等式失败。平方误差放大了大的差异，经由中间点分两步走反而“更近”；这正是它不能直接当作距离使用的原因，而它的平方根$\lvert x-y\rvert$是度量。

\subsection{度量拓扑}
\label{sec:metric-topology}

有了距离，“$x$附近”最自然的含义是距离$x$小于某个正数的那些点。我们先给这类集合命名，再说明它们怎样生成第~\ref{chap:topological-spaces}章意义下的拓扑。

\begin{definition}[开球与闭球]
\label{def:metric-balls}
设$(X,d)$是度量空间，$x\in X$。对$r>0$，称
\begin{equation*}
B(x,r)=\{y\in X:d(x,y)<r\},\qquad B[x,r]=\{y\in X:d(x,y)\leq r\}
\end{equation*}
分别为以$x$为中心、$r$为半径的\term{开球}（open ball）与\term{闭球}（closed ball）。需要指明度量时写作$B_d(x,r)$。
\end{definition}

方括号提示闭球包含距离恰为$r$的点。闭球不一定是开球的闭包：在离散度量$d_0$下，$B(x,1)=\{x\}$，它的闭包仍是$\{x\}$，而$B[x,1]$是整个空间。所以本书不用$\overline{B(x,r)}$表示闭球。

开球的关键性质是：球内每一点都被一个更小的球包在球内。

\begin{lemma}[开球内的小球]
\label{lem:metric-ball-inside-ball}
若$y\in B(x,r)$，令$s=r-d(x,y)>0$，则$B(y,s)\subseteq B(x,r)$。
\end{lemma}

\begin{proof}
对$z\in B(y,s)$，由三角不等式，$d(x,z)\leq d(x,y)+d(y,z)<d(x,y)+s=r$，所以$z\in B(x,r)$。
\end{proof}

\begin{theorem}[开球构成拓扑基]
\label{thm:metric-balls-base}
度量空间$(X,d)$中全体开球构成$X$上某个拓扑的基。由这个基生成的拓扑中，$U\subseteq X$是开集当且仅当对每个$x\in U$都存在$r>0$使$B(x,r)\subseteq U$。
\end{theorem}

\begin{proof}
按第~\ref{chap:topological-spaces}章拓扑基的判定，要核对两件事。第一，每个$x$都属于开球$B(x,1)$，所以开球覆盖$X$。第二，设$z\in B(x_1,r_1)\cap B(x_2,r_2)$，由引理~\ref{lem:metric-ball-inside-ball}，取$s=\min\{r_1-d(x_1,z),r_2-d(x_2,z)\}>0$，则$B(z,s)$同时含于两个球，因而含于它们的交。

由基生成的开集是若干基元素的并。若$U$是开球之并，$x\in U$落在某个$B(y,t)\subseteq U$中，引理~\ref{lem:metric-ball-inside-ball}给出以$x$为中心、含于$B(y,t)$的球。反过来，若每个$x\in U$都有$B(x,r_x)\subseteq U$，则$U=\bigcup_{x\in U}B(x,r_x)$是开球之并。
\end{proof}

由开球生成的拓扑称为$d$的\term{度量拓扑}（metric topology）。今后说度量空间中的开集、闭集、邻域、收敛与连续，都指这个拓扑下第~\ref{chap:topological-spaces}章的概念。按定理~\ref{thm:metric-balls-base}，集合$N$是$x$的邻域当且仅当它包含某个开球$B(x,r)$；于是$B(x,1/n)$，$n=1,2,\ldots$，构成$x$的一个可数邻域基，因为任何$B(x,r)$都包含$n>1/r$时的$B(x,1/n)$。

\paragraph{点的分离}
距离为正的两点可以用半径为距离一半的球分开。设$x\neq y$，$r=d(x,y)/2>0$。若$z\in B(x,r)\cap B(y,r)$，则$d(x,y)\leq d(x,z)+d(z,y)<2r=d(x,y)$，矛盾，所以$B(x,r)$与$B(y,r)$不交。这说明度量拓扑满足Hausdorff条件：不同的点有不交的邻域。

这条性质给出一个限制：并非每个拓扑都来自某个度量。两点集$\{a,b\}$上的平凡拓扑只有$\varnothing$和全集两个开集，$a,b$无法用不交的开集分开，所以任何度量都不能诱导它。另一个直接推论是单点集是闭集：对$y\neq x$，球$B(y,d(x,y))$不含$x$，所以$X\setminus\{x\}$是开集。

\paragraph{用球判断子集的位置}
把第~\ref{chap:topological-spaces}章的定义换成球的语言，子集$A\subseteq X$与点$x$的关系可以这样判断：$x$是$A$的内点，当且仅当某个$B(x,r)\subseteq A$；$x$属于$A$的闭包，当且仅当每个$B(x,r)$都与$A$相交；$x$属于$A$的边界，当且仅当每个$B(x,r)$既与$A$相交、又与$X\setminus A$相交。这三条都由“开球构成邻域基”直接得到。

闭球是闭集。若$d(x,y)>r$，令$s=d(x,y)-r$，对$z\in B(y,s)$有$d(x,z)\geq d(x,y)-d(y,z)>r$，所以$B(y,s)$落在$B[x,r]$的补集中，补集是开集。

\begin{example}[开闭性与边界的判断]
\label{ex:metric-open-closed-interval}
在$\mathbb R$的通常度量$\lvert x-y\rvert$下考察$A=[0,1)$。点$0$的每个球都含有负数，所以$0$不是内点，$A$不是开集；点$1$的每个球都含有$A$中的点，$1$在闭包中却不在$A$中，所以$A$也不是闭集。按上面的判据，$A$的内部为$(0,1)$，闭包为$[0,1]$，边界为$\{0,1\}$。

换成离散度量$d_0$，每个单点集都是开球$B(x,1/2)$，于是任意子集都是开集，它的补集也是开集，所以任意子集同时是开集和闭集，边界总是空集。同一个集合$[0,1)$在两种度量下开闭性不同，再次说明开闭性取决于所选的度量。
\end{example}

不同的度量也可能给出相同的拓扑。在$\mathbb R^2$中，对任意$\bm x$和$r>0$，
\begin{equation*}
B_{d_\infty}(\bm x,r/\sqrt2)\subseteq B_{d_2}(\bm x,r)\subseteq B_{d_\infty}(\bm x,r),
\end{equation*}
这由$d_\infty\leq d_2\leq\sqrt2\,d_\infty$得到。于是一种度量下的每个球内都含有另一种度量下同心的球，按定理~\ref{thm:metric-balls-base}，两者的开集完全相同。离散度量$d_0$则给出不同的拓扑：单点集在$d_0$下是开集，在$d_2$下不是。什么时候两个度量“一样好”，取决于要保留开集、柯西列还是误差尺度，第~\ref{chap:metric-space-maps}章分别讨论。

\subsection{有界集}
\label{sec:metric-bounded}

拓扑只告诉我们哪些点“附近”，不告诉我们一个集合整体有多大。距离可以回答后一个问题：一个集合能否被某个有限半径的球装下。

\begin{definition}[有界集与直径]
\label{def:metric-bounded}
度量空间$(X,d)$的子集$A$若满足：存在$x\in X$和$r>0$使$A\subseteq B(x,r)$，则称$A$为\term{有界集}（bounded set）。非空集合$A$的\term{直径}（diameter）定义为$\operatorname{diam}A=\sup\{d(a,b):a,b\in A\}\in[0,\infty]$。空集规定为有界。
\end{definition}

球心的选取并不重要。若$A\subseteq B(x,r)$，对任意$y\in X$和$a\in A$，有$d(y,a)\leq d(y,x)+d(x,a)<d(y,x)+r$，所以$A\subseteq B(y,r+d(x,y))$：换一个球心，只需把半径加上两个球心的距离。同样的估计说明，非空集合有界当且仅当它的直径有限：若$A\subseteq B(x,r)$，则$d(a,b)\leq d(a,x)+d(x,b)<2r$；反过来，若$\operatorname{diam}A<\infty$，任取$a_0\in A$，就有$A\subseteq B(a_0,\operatorname{diam}A+1)$。

有界性与开闭性互不决定。在$\mathbb R$中，$(0,1)$有界且是开集；整数集$\mathbb Z$是闭集，因为它的补集是开区间之并，但它不能被任何有限半径的球容纳；$[0,1]$既有界又闭。可见判断一个集合“大小有限”，不能从它是否包含边界点来推断。

有界性还依赖于度量，而不只依赖于拓扑。在度量空间$(X,d)$上令$d'(x,y)=\min\{d(x,y),1\}$。前两条公理显然成立；对三角不等式，若$d(x,y)\geq1$或$d(y,z)\geq1$，右端至少为$1$，不小于左端；否则右端等于$d(x,y)+d(y,z)\geq d(x,z)\geq d'(x,z)$。半径小于$1$的球在$d$和$d'$下完全相同，所以两者拓扑相同，但在$d'$下每个集合都含于$B_{d'}(x,2)$，整个$\mathbb R$也成了有界集。

\subsection{可分性的判据}
\label{sec:metric-separability}

第~\ref{chap:topological-spaces}章区分了两种可数性：局部上，每个点有可数邻域基；整体上，整个空间有可数拓扑基。还定义了可分性：存在可数的稠密子集。在度量空间中，前一种总是成立，后两种恰好等价。

局部可数性已经在第~\ref{sec:metric-topology}小节得到：$\{B(x,1/n)\}_{n\geq1}$是$x$的邻域基。但这不能推出整个空间有可数拓扑基，因为每个点各自的可数个球合起来，在点数不可数时就不再可数。$\mathbb R$配上离散度量$d_0$就是例子：每个单点集$\{x\}=B(x,1/2)$都是开集，任何拓扑基必须为每个$x$提供一个含$x$且含于$\{x\}$的基元素，即$\{x\}$本身，所以拓扑基至少有不可数个元素。

\begin{theorem}[度量空间的可分性]
\label{thm:metric-separable-second-countable}
度量空间$(X,d)$可分，当且仅当它有可数拓扑基。
\end{theorem}

\begin{proof}
先设$D\subseteq X$可数且稠密。考虑族$\mathcal B=\{B(q,r):q\in D,\ r\in\mathbb Q,\ r>0\}$，它由两个可数集的乘积编号，所以可数。设$U$是开集，$x\in U$，取$r>0$使$B(x,r)\subseteq U$，再取有理数$s$满足$0<s<r/2$。由稠密性，开球$B(x,s)$含有某个$q\in D$。于是$x\in B(q,s)$，并且对$z\in B(q,s)$有$d(x,z)\leq d(x,q)+d(q,z)<2s<r$，即$B(q,s)\subseteq B(x,r)\subseteq U$。所以每个开集都是$\mathcal B$中元素之并，$\mathcal B$是拓扑基。

反过来，设$\{U_k\}_{k\geq1}$是可数拓扑基。对每个非空的$U_k$取一点$x_k\in U_k$，令$D$为这些点组成的集合，它可数。任意非空开集$V$是基元素之并，所以包含某个非空的$U_k$，从而含有$x_k\in D$。每个非空开集都与$D$相交，$D$稠密。
\end{proof}

证明的第一步用有理半径，是为了让球的总数可数；第二步只是从每个基元素中选一个代表点。在一般拓扑空间中，第二步仍然成立，即有可数基的空间可分；第一步却用到了距离和三角不等式，一般拓扑空间中可分不保证有可数基。

\begin{example}[欧氏空间的可数基]
\label{ex:metric-rn-separable}
在$(\mathbb R^n,d_2)$中，有理坐标点集$\mathbb Q^n$可数。给定$\bm x$和$\varepsilon>0$，逐个坐标取有理数$q_i$使$\lvert x_i-q_i\rvert<\varepsilon/\sqrt n$，则$d_2(\bm x,\bm q)<\varepsilon$，所以$\mathbb Q^n$稠密。按定理~\ref{thm:metric-separable-second-countable}，以有理坐标点为中心、有理数为半径的球构成可数拓扑基。前面的离散$\mathbb R$没有可数基，相应地它也不可分：每个单点集都是开集，稠密子集必须与每个单点集相交，只能是$\mathbb R$本身。
\end{example}

这个等价关系还有一个推论：可分度量空间的每个子集，按限制下来的度量也可分。原因是可数拓扑基$\{U_k\}$与子集$A$的交$\{U_k\cap A\}$构成$A$的可数拓扑基；这里用到子集上的拓扑恰是限制度量给出的拓扑，第~\ref{sec:metric-subspace}小节会证明这一点。例如$\mathbb R^2$中的任意曲线都可分，推论直接给出这一点，不需要具体找出它的可数稠密子集；对无理数集，也可以直接核对$\{q+\sqrt2:q\in\mathbb Q\}$是它的可数稠密子集。
